\section{Dataset Overview}

\subsection{Characteristics of \benchmarkname}

\textbf{Industrially-Relevant, Diverse, and Challenging Tasks.} First, all repositories selected in \benchmarkname~are actively maintained professional projects with substantial user bases, comprehensive documentation, and established development practices. In addition, we source commercial repositories. These repositories are private and sourced from startups, where we contacted the company and purchased their engineering repos. We sample repositories from a diverse range of topics, including consumer applications with complex UI logic, B2B platforms with intricate business rules, and developer tools with sophisticated APIs. Second, we limit each repository to contribute 50-100+ instances. This avoids the situation where models get an advantage by being especially good at a single repository, rewarding models that can truly generalize. Finally, we require edits to span multiple files and contain a substantial code change, similar to real software engineering tasks. Subsequently, \benchmarkname~problems are naturally challenging.


\textbf{Verified and Human-Augmented.} Similar to SWE-Bench Verified, each problem in \benchmarkname~goes through a human augmentation and verification process. This ensures that task descriptions are not missing critical information, tests are well specified to validate the generated solution, and problems are representative of real-world software engineering tasks. In particular, we augment each issue with a list of human-written requirements -- simulating the standard engineering practice of resolving issues follow problem specification and provide additional guarantee that the problems are self-contained. Note that real software engineering tasks can be under-specified (for example, may require exploration before solving), and that the setting \textit{without requirements} is potentially interesting.

\textbf{Contamination-Resistant by Design.} By exclusively using repositories with GPL and other copyleft licenses, we ensure benchmark content is unlikely to appear in proprietary model training sets, as the nature of these licenses creates legal barriers to their inclusion in commercial training corpora. In addition, we use commercial repositories purchased from startups, which are private.

\subsection{Task Specification}

% Each task instance in \benchmarkname~is complete with a problem statement, requirements, interfaces, and a complete environment for the agent. Below, we give detail on the task specification in \benchmarkname. A model receives three components: (1) a human-augmented problem statement providing comprehensive context about the problem, including current behavior, expected behavior, and relevant background information; (2) an explicit enumeration of requirements that must be satisfied for successful resolution, formatted as testable conditions; and (3) a complete repository codebase from a professional software project. The model must generate a patch file specifying modifications to resolve the issue while meeting all requirements.

Each task instance in \benchmarkname~is complete with human-augmented problem statement, requirements and interface as the task description for the model. The model must generate a patch file to resolve the issue and pass a suite of human-reviewed tests as validation.

% \textbf{Problem Statement.} Similar to SWE-Bench \cite{zhang2025swebench}, we provide a problem statement containing a task definition, similar to the description of an issue or pull-request in a Github repository. Agents should be able to solve the task using only the problem statement.

\textbf{Problem Statement.} Similar to SWE-Bench, we provide a problem statement describing the issue to solve. We use content from the original commits, PR and issue, then rewrite it in the style of issues and add in missing information when necessary. Agents should be able to solve the task using only the problem statement.

\textbf{Requirements.} Problems in \benchmarkname~can be more complex than previous iterations of SWE-Bench, and thus, we introduce requirements to resolve any potential ambiguity issues. For each problem, we list out a set of requirements that give additional detail on what is needed to solve the task. These requirements are grounded on the unit tests that are used for validation. For example, a requirement might specify the route names and functionality expected for an API.

\textbf{Interface.}
\benchmarkname~tasks include feature additions and code enhancements, such as refactoring, that involve creating or altering classes and methods. For these tasks, a common false negative in unit-test verifiers is when a model submits a valid solution with different interfaces than what the unit test is expecting. Here, we explicitly define the class and function names expected by the tests to avoid this failure mode when relevant.


% {models may choose class or function names that mismatch with what is used in tests as the interface is not clearly specified in the problem statement.} 

\textbf{Environments.} Each task is evaluated in a containerized, language-specific environment with full dependency resolution. Python tasks use isolated virtual environments, JavaScript/TypeScript tasks use Node.js with npm/yarn, and Go tasks use module-aware environments with proper GOPATH configuration. All environments will be released as pre-built docker images to ensure that they are fully reproducible.

\textbf{Tests.} Every task includes human-reviewed test suites with \texttt{fail2pass} tests that verify issue resolution and \texttt{pass2pass} tests that ensure existing functionalities remain intact. % We first run the tests without the gold patch, then apply the gold patch to determine relevant test statuses. We notice that some tests can be dynamic or fail occasionally. To mitigate it, we run each set of tests 3 times and filter out any test that doesn't pass consistently. Finally, we perform an additional round of verification on the \texttt{fail2pass} tests where we ask annotators to filter out tests which are too broad or not relevant to the task description.

\subsection{Public, Commercial, and Held-Out \benchmarkname}

\benchmarkname~consists of a total of 1865 human-verified and augmented problems, divided as three subsets: public, commercial, and held-out.
\begin{itemize}
\item \textbf{Public}. We release 731 instances openly on HuggingFace and report the relevant statistics and model performances in this paper. These are sourced from public repositories with copy-left license. 
\item \textbf{Commercial}. For the commercial set of 276 problems sourced from startup repositories, we keep it private but report results publicly in this paper and will update in the leaderboard. This is the only set containing proprietary repositories from startups, which we cannot release for legal reasons.
\item \textbf{Held-Out}. We hold out a set of 858 problems mirroring the public set but use a separate set of repositories. We keep this set private to test for overfitting in the future.
\end{itemize}